\documentclass[11pt]{article}
\usepackage[margin=1in]{geometry}
\usepackage{amsmath, amssymb, amsthm}
\usepackage{mathtools}
\usepackage{hyperref}
\usepackage{algorithm}
\usepackage{algpseudocode}
\usepackage{graphicx}
\usepackage{float}
\usepackage[numbers,sort&compress]{natbib}

\newtheorem{proposition}{Proposition}
\newtheorem{definition}{Definition}

\title{What Algorithms Can Transformers Learn In-Context?\\ Online selection of online algorithms}
\author{Senem I\c{s}\i{}k}

\begin{document}
\maketitle

\tableofcontents
\bigskip

\section{Introduction}
\label{sec:intro}

In-context learning is the phenomenon by which a sequence model, at test time, executes a procedure inferred from the data it sees in context --- potentially noisy input-output pairs only --- without any weight updates. Prior work has demonstrated this for simple function classes such as linear and boolean functions, but not for online algorithms. Yet many real decision problems are solved not by fitting a function but by running such an algorithm: integrating the data so far, updating beliefs, choosing an adaptive strategy, and acting. Can transformers learn this kind of instance adaptability --- recovering and running the right algorithmic strategy for each evolving sequence --- or do they collapse onto a single fixed procedure tuned to the average of the training distribution?

Our testbed is a Bayesian stopping problem with a Normal--Normal conjugate prior, in which the optimal policy is a backward-induction threshold rule indexed by a latent noise scale $\sigma$. We call this setting \emph{online selection of online algorithms}: at test time the transformer sees a sequence with unknown noise level and must, on the fly, recover from the prefix it has observed the threshold rule that corresponds to that noise level.

By varying whether the noise level is fixed or random across training distributions, we engineer a setting that discriminates between two hypotheses about the trained model:
\begin{itemize}
\item \textbf{Learning-algorithm hypothesis.} The transformer infers the noise level in-context and applies the corresponding threshold rule, adapting per instance and generalizing across noise levels.
\item \textbf{Shortcut hypothesis.} The transformer applies a single training-distribution-tuned policy that ignores the noise level, matching the oracle in-distribution but degrading sharply elsewhere.
\end{itemize}
To tell them apart, we analyze the trained network with an interpretability stack --- attention entropy, layerwise probing, residual-stream interventions, activation patching, head ablation, and a logit lens through depth --- asking which intermediate quantities the model computes, where in the network they live, and whether they causally drive the output. The goal is to support the case that transformers can select online algorithms in-context (with the right training), and to explain how.

Section~\ref{sec:results} reports the empirical findings. Multi-regime training under $\mathcal{D}_{\mathrm{disc}}$ confirms the algorithm hypothesis. Training under the continuous prior $\mathcal{D}_{\mathrm{logu}}$ surfaces a supervision-form-dependent failure --- continuation-value supervision degrades at low $\sigma$ while action supervision recovers fully --- which the diagnostic in Section~\ref{sec:res-logu-cv} localizes to the linear continuation-value head.

\section{Setup: A Bayesian Optimal Stopping Problem}
\label{sec:setup}

\subsection{Generative model}

For each sample $i \in \{1, \ldots, N\}$, draw a within-sequence noise scale $\sigma_i$ from a prior, draw a latent mean $\mu_i \sim \mathcal{N}(\mu_0, \tau_0^2)$, then draw a sequence
\[
X_1, X_2, \ldots, X_n \mid \mu_i, \sigma_i \;\overset{\text{i.i.d.}}{\sim}\; \mathcal{N}(\mu_i, \sigma_i^2).
\]
We fix $\mu_0 = 0$, $\tau_0 = 10$, horizon $n = 256$ throughout. The agent observes $X_t$ at each $t \in \{1, \ldots, n\}$ and chooses to \emph{stop} (collect $X_t$) or \emph{continue}; at $t = n$ it must stop. Both $\mu_i$ and (in the unknown-$\sigma$ setting) $\sigma_i$ are latent.

\subsection{Five training distributions}

We instantiate five training distributions, three with $\sigma$ fixed (\emph{static-variance}) and two with $\sigma$ random (\emph{random-variance}):
\begin{itemize}
\item $\mathcal{D}_1$: static, $\sigma = 1$ (data-dominant).
\item $\mathcal{D}_2$: static, $\sigma = 10$ (balanced).
\item $\mathcal{D}_3$: static, $\sigma = 100$ (prior-dominant).
\item $\mathcal{D}_{\mathrm{disc}}$: random, $\sigma \sim \mathrm{Uniform}\{1, 10, 100\}$, each with probability $1/3$.
\item $\mathcal{D}_{\mathrm{logu}}$: random, $\log_{10}\sigma \sim \mathrm{Uniform}(0, 2)$, equivalently $p(\sigma) = 1/(2\sigma\ln 10)$ on $[1, 100]$.
\end{itemize}
Both random-variance priors give equal weight to each decade of $\sigma$ --- the discrete prior trivially, the continuous prior by uniform-in-$\log_{10}\sigma$ construction.

The static-variance distributions $\mathcal{D}_1, \mathcal{D}_2, \mathcal{D}_3$ play a double role: training distributions for the single-regime models, and test regimes for all evaluation (Section~\ref{sec:regimes}). We name single-regime models by their training data (D1\_cv, D1\_act, etc.) when we need to disambiguate.

\subsection{Three test regimes}
\label{sec:regimes}

The static-variance distributions $\mathcal{D}_1, \mathcal{D}_2, \mathcal{D}_3$ are also the three test regimes for all evaluation. They correspond to noise-to-prior-precision ratios $\rho := \sigma^2/\tau_0^2$, with the data weight at $t = 1$ given by $1/(1 + \rho)$:

\begin{center}
\small
\begin{tabular}{l c c c c l}
\hline
test regime & $\sigma$ & $\sigma^2$ & $\rho$ & data weight at $t = 1$ & character \\
\hline
$\mathcal{D}_1$ (data-dominant)    & $1$   & $1$     & $0.01$ & $0.99$ & plug-in is asymptotically Bayes \\
$\mathcal{D}_2$ (balanced)         & $10$  & $100$   & $1$    & $0.50$ & no closed-form is asymptotic \\
$\mathcal{D}_3$ (prior-dominant)   & $100$ & $10000$ & $100$  & $0.01$ & prior-only is asymptotically Bayes \\
\hline
\end{tabular}
\end{center}

\paragraph{Datasets.} Training sequences are sampled fresh from the training distribution at each step (no fixed training set); the meaningful budget is gradient steps $\times$ batch size, specified in Section~\ref{sec:training}. Validation and test sets are pre-generated and held fixed at $N_{\mathrm{val}} = N_{\mathrm{test}} = 10^4$ sequences each (validation drawn from the training distribution, used for checkpoint selection only; test drawn per regime $\mathcal{D}_i$, used for all evaluation). Validation and test sets use disjoint random seeds.

\section{Bayes-optimal Oracle}
\label{sec:oracle}

\subsection{Static-variance setting}
\label{sec:oracle-static}

When the training distribution is one of $\mathcal{D}_1, \mathcal{D}_2, \mathcal{D}_3$ the noise scale $\sigma$ is constant and known; only $\mu$ is latent. The oracle is the standard Normal--Normal stopping DP.

\paragraph{Sufficient statistic and posterior.} The sufficient statistic for $\mu$ given $X_{1:t}$ is the running sum
\[
S_t \;:=\; \sum_{s=1}^t X_s,
\]
with derived sample mean $\bar X_t := S_t / t$ (the maximum-likelihood estimate of $\mu$ at known $\sigma$). The conjugate Normal--Normal posterior depends on $X_{1:t}$ only through $S_t$:
\begin{equation}
\label{eq:mu-static}
\mu \mid X_{1:t} \;\sim\; \mathcal{N}(\mu_t, \tau_t^2),
\qquad
\tau_t^2 = \!\left(\frac{1}{\tau_0^2} + \frac{t}{\sigma^2}\right)^{\!-1},
\qquad
\mu_t = \tau_t^2\!\left(\frac{\mu_0}{\tau_0^2} + \frac{S_t}{\sigma^2}\right).
\end{equation}
$\tau_t^2$ is deterministic in $t$; $\mu_t$ is a deterministic function of $(t, S_t)$. The belief state is $b_t = \mu_t$, one-dimensional.

\paragraph{DP recursion.} The continuation value satisfies $V_t(b_t, X_t) = \max\{X_t, C_t(\mu_t)\}$ for $t < n$ and $V_n = X_n$, with
\begin{equation}
\label{eq:Crec-static}
C_t(\mu_t) \;=\; \mathbb{E}_{X \sim \mathcal{N}(\mu_t,\, \tau_t^2 + \sigma^2)}\!\Big[\max\!\big\{X,\, C_{t+1}(\mathrm{update}(\mu_t, X))\big\}\Big],
\end{equation}
where $\mathrm{update}(\mu, X) = \alpha_t \mu + (1 - \alpha_t) X$ with $\alpha_t := \sigma^2/(\tau_t^2 + \sigma^2)$, and $C_{n-1}(\mu) = \mu$. The Bayes-optimal policy is $\pi^\star_{\mathcal{D}_i, t}(b_t, X_t) = \mathbf{1}[X_t \ge C_t(\mu_t)]$.

\paragraph{Approximate DP.}
\label{sec:adp-static}
Equation~\eqref{eq:Crec-static} has no elementary closed form. We discretize $\mu_t$ on a uniform grid of $K$ points covering $[\mu_0 - 5 s_t,\, \mu_0 + 5 s_t]$ at each stage, where $s_t = \sqrt{\tau_0^2 - \tau_t^2}$ is the marginal standard deviation of $\mu_t$ across sequences. The Gaussian expectation in~\eqref{eq:Crec-static} is approximated by $J$-point Gauss--Hermite quadrature, the standard rule for integrals against $e^{-z^2}$. The nodes $\{z_j\}_{j=1}^J \subset \mathbb{R}$ are the roots of the $J$-th physicists' Hermite polynomial $H_J$, and the weights $\{w_j\}_{j=1}^J > 0$ are uniquely determined by requiring that
\[
\int_{-\infty}^\infty f(z)\, e^{-z^2}\, dz \;\approx\; \sum_{j=1}^J w_j\, f(z_j)
\]
hold with equality whenever $f$ is a polynomial of degree at most $2J - 1$. Both nodes and weights depend only on $J$ and are precomputed (e.g., via \texttt{numpy.polynomial.hermite.hermgauss}). Substituting $X = \mu + \sqrt{2 v_t}\,z$ for $X \sim \mathcal{N}(\mu, v_t)$ with $v_t := \tau_t^2 + \sigma^2$ gives
\begin{equation}
\label{eq:gh}
\mathbb{E}_{X \sim \mathcal{N}(\mu, v_t)}[f(X)] \;\approx\; \frac{1}{\sqrt\pi}\sum_{j=1}^J w_j\, f\!\left(\mu + \sqrt{2 v_t}\, z_j\right).
\end{equation}
Off-grid values of $\widehat C_{t+1}$ are looked up by linear interpolation with clipping at the grid boundary. Algorithm~\ref{alg:adp-static} runs the backward induction at $K = 2048$, $J = 128$.

\begin{algorithm}[H]
\caption{Approximate DP for the static-variance Bayes-optimal threshold (per regime).}
\label{alg:adp-static}
\begin{algorithmic}[1]
\State Precompute $\tau_t^2, v_t := \tau_t^2 + \sigma^2, \alpha_t$ for each $t$; grid $\{m_t^{(k)}\}_{k=1}^K$ at each $t$; Gauss--Hermite nodes $\{(z_j, w_j)\}_{j=1}^J$
\State $\widehat C_{n-1}(m_{n-1}^{(k)}) \gets m_{n-1}^{(k)}$ for $k = 1, \ldots, K$
\For{$t = n-2, n-3, \ldots, 1$}
  \For{$k = 1, \ldots, K$}
    \State $S \gets 0$
    \For{$j = 1, \ldots, J$}
      \State $x_j \gets m_t^{(k)} + \sqrt{2 v_t}\, z_j$
      \State $\mu_j' \gets \alpha_t\, m_t^{(k)} + (1-\alpha_t)\, x_j$
      \State $S \gets S + w_j\, \max\!\big\{x_j,\, \widehat C_{t+1}^{\mathrm{lin}}(\mu_j')\big\}$
    \EndFor
    \State $\widehat C_t(m_t^{(k)}) \gets S/\sqrt\pi$
  \EndFor
\EndFor
\State \Return $\{\widehat C_t\}$
\end{algorithmic}
\end{algorithm}

\paragraph{Convergence check.} Solve at $(K, J) = (2048, 128)$ and at reference $(K, J) = (4096, 256)$; report (i) max absolute disagreement on $\widehat C$ over the grid, and (ii) action-label disagreement on $10^4$ test sequences per regime. Fix the labeling resolution after the operational disagreement falls below $10^{-3}$.

\subsection{Random-variance setting}
\label{sec:oracle-random}

When the training distribution is $\mathcal{D}_{\mathrm{disc}}$ or $\mathcal{D}_{\mathrm{logu}}$, both $\mu$ and $\sigma$ are latent. The oracle is a two-dimensional DP over a joint posterior on $(\mu, \sigma)$.

\paragraph{Sufficient statistics.} The sufficient statistics for the joint posterior are
\begin{equation}
\label{eq:S-Q-defs}
S_t \;:=\; \sum_{s=1}^t X_s,
\qquad
Q_t \;:=\; \sum_{s=1}^t X_s^2.
\end{equation}
Both update online: $S_t = S_{t-1} + X_t$, $Q_t = Q_{t-1} + X_t^2$. Together with $t$, the pair $(S_t, Q_t)$ determines the joint posterior $p(\mu, \sigma \mid X_{1:t})$. $S_t$ alone determines the sample mean $\bar X_t = S_t/t$; jointly $(S_t, Q_t)$ determine the within-sequence sample variance
\[
\hat\sigma_t^2 \;:=\; \frac{1}{t-1}(Q_t - t\bar X_t^2).
\]

\paragraph{Joint posterior.} By Bayes' rule, $p(\mu, \sigma \mid X_{1:t}) \propto p(\sigma)\, p(\mu)\, p(X_{1:t} \mid \mu, \sigma)$. Factor as $p(\sigma \mid X_{1:t}) \cdot p(\mu \mid X_{1:t}, \sigma)$. For any fixed $\sigma$, the conditional posterior $p(\mu \mid X_{1:t}, \sigma)$ is the standard Normal--Normal posterior~\eqref{eq:mu-static}: depends on $X_{1:t}$ only through $S_t$, with different $\sigma$ values giving different conditional means and variances.

For the marginal posterior $p(\sigma \mid X_{1:t})$: the marginal likelihood $p(X_{1:t} \mid \sigma) = \int p(X_{1:t} \mid \mu, \sigma)\, p(\mu)\, d\mu$ has $X_{1:t} \mid \sigma$ multivariate Gaussian with mean $\mathbf{0}$ and covariance $\Sigma_t^{(\sigma)} = \sigma^2 I_t + \tau_0^2 \mathbf{1}\mathbf{1}^\top$. Using the matrix determinant lemma and Sherman--Morrison (Appendix~\ref{app:marginal-derivation}), the log marginal likelihood depends on $X_{1:t}$ only through $(S_t, Q_t)$:
\begin{equation}
\label{eq:marginal-loglik}
\log p(X_{1:t} \mid \sigma) = \mathrm{const}(t) - \frac{t}{2}\log\sigma^2 - \frac{1}{2}\log\!\left(1 + \frac{t\tau_0^2}{\sigma^2}\right) - \frac{1}{2}\!\left[\frac{Q_t}{\sigma^2} - \frac{\tau_0^2 S_t^2/\sigma^4}{1 + t\tau_0^2/\sigma^2}\right].
\end{equation}
For $\mathcal{D}_{\mathrm{disc}}$ the posterior over $\sigma$ is a categorical on three points; for $\mathcal{D}_{\mathrm{logu}}$ a continuous density on $[1, 100]$. Evaluation details are deferred to the Approximate DP paragraph below.

\paragraph{Predictive distribution and DP recursion.} The predictive $p(X_{t+1} \mid X_{1:t})$ is a mixture of Gaussians with weights from $p(\sigma \mid X_{1:t})$:
\begin{equation}
\label{eq:predictive}
p(X_{t+1} \mid X_{1:t}) \;=\; \int p(\sigma \mid X_{1:t}) \cdot \mathcal{N}\!\left(X_{t+1};\, \mu_t^{(\sigma)},\, v_t^{(\sigma)}\right)\, d\sigma,
\qquad v_t^{(\sigma)} := \tau_t^{2,(\sigma)} + \sigma^2,
\end{equation}
with the integral replaced by a sum over $\{1, 10, 100\}$ for $\mathcal{D}_{\mathrm{disc}}$. The belief state is $b_t = (S_t, Q_t)$. Backward induction gives $V_t(b_t, X_t) = \max\{X_t, C_t(b_t)\}$ for $t < n$ with $V_n = X_n$, and
\begin{equation}
\label{eq:Crec-random}
C_t(b_t) \;=\; \mathbb{E}_{X_{t+1} \sim p(\cdot \mid b_t)}\!\big[\max\{X_{t+1}, C_{t+1}(b_{t+1})\}\big],
\end{equation}
expectation over the mixture predictive~\eqref{eq:predictive}, $b_{t+1} = (S_t + X_{t+1}, Q_t + X_{t+1}^2)$, terminal $C_{n-1}(b_{n-1}) = \int p(\sigma \mid X_{1:n-1}) \mu_{n-1}^{(\sigma)}\, d\sigma$. The Bayes-optimal policy is $\pi^\star_{\mathcal{D}, t}(b_t, X_t) = \mathbf{1}[X_t \ge C_t(b_t)]$.

\paragraph{Approximate DP.}
\label{sec:adp-random}
We discretize $b_t$ in transformed coordinates $(\bar X_t, L_t)$ where $L_t := \log\hat\sigma_t^2$. The two axes need different treatment:
\begin{itemize}
\item $\bar X_t = \mu + (1/t)\sum_s \varepsilon_s$ has marginal variance $\tau_0^2 + \sigma^2/t$ across sequences; at fixed $\sigma$, this only converges to $\tau_0^2$ as $t \to \infty$. At small $t$ and large $\sigma$ the marginal SD is much larger than $\tau_0$, so a fixed grid is unsuitable. We instead use a \emph{per-stage} uniform grid $I_{\bar X, t} := [-5\, s_{\bar X, t},\, +5\, s_{\bar X, t}]$ with stage-dependent half-width $s_{\bar X, t} := \sqrt{\tau_0^2 + \sigma_{\max}^2 / t}$, where $\sigma_{\max}$ is the supremum of the prior's support ($\sigma_{\max} = 100$ for both $\mathcal{D}_{\mathrm{disc}}$ and $\mathcal{D}_{\mathrm{logu}}$). At $t = 1$: $s_{\bar X, 1} = \sqrt{\tau_0^2 + \sigma_{\max}^2} \approx 100.5$, grid spans $\approx [-503, 503]$. As $t \to \infty$: $s_{\bar X, t} \to \tau_0 = 10$, grid spans $\approx [-50, 50]$. Mirrors the static ADP's per-stage $\mu_t$ grid (Section~\ref{sec:adp-static}).
\item $L_t$ estimates $\log\sigma^2$, which spans $[\log 1, \log 10^4] = [0, 4\log 10]$. Grid on $I_L := [\log 0.1, \log 10^5]$ to absorb sample-variance dispersion at small $t$. Stage-invariant.
\end{itemize}
Use $K_1 = K_2 = 256$ uniform grid points on each axis, total $K_1 K_2 = 65{,}536$ states per stage. The state translation between $(S_t, Q_t)$ and grid coordinates is $S_t = t \bar X_t$, $Q_t = (t-1) e^{L_t} + t \bar X_t^2$.

For the mixture quadrature, use $M$ components: $M = 3$ for $\mathcal{D}_{\mathrm{disc}}$ at $\sigma \in \{1, 10, 100\}$, or $M = J_\sigma = 64$ Gauss--Legendre nodes on $\log\sigma \in [0, \log 100]$ for $\mathcal{D}_{\mathrm{logu}}$ (Gauss--Legendre is the analog of Gauss--Hermite for integrals against the uniform measure on a finite interval). For each component $\sigma_k$ apply $J = 64$-point Gauss--Hermite~\eqref{eq:gh}; sum the per-component results weighted by $\tilde w_k := p(\sigma_k \mid X_{1:t}) \cdot \omega_k / Z_t$, where $\omega_k$ is the Gauss--Legendre weight (or $1$ for the discrete case) and $Z_t$ normalizes. Off-grid lookup uses 2D bilinear interpolation in $(\bar X_{t+1}, L_{t+1})$.

\begin{algorithm}[H]
\caption{Approximate DP for the random-variance Bayes-optimal threshold.}
\label{alg:adp-random}
\begin{algorithmic}[1]
\State \textbf{Precompute.} Per-stage $\bar X$ grids $\{\bar X_t^{(k_1)}\}$ on $I_{\bar X, t} = [-5 s_{\bar X, t}, +5 s_{\bar X, t}]$ with $s_{\bar X, t} = \sqrt{\tau_0^2 + \sigma_{\max}^2 / t}$; stage-invariant $L$ grid $\{L^{(k_2)}\}$ on $I_L$.
\State \textbf{Initialize.} For each grid point $(k_1, k_2)$, recover $(S_{n-1}, Q_{n-1})$ from $(t = n-1, \bar X_{n-1}^{(k_1)}, L^{(k_2)})$, compute $p(\sigma \mid X_{1:n-1})$ via~\eqref{eq:marginal-loglik}, set $\widehat C_{n-1}(\bar X_{n-1}^{(k_1)}, L^{(k_2)}) \gets \int p(\sigma \mid X_{1:n-1})\, \mu_{n-1}^{(\sigma)}\, d\sigma$
\For{$t = n-2, n-3, \ldots, 1$}
  \For{each grid point $(k_1, k_2)$ on stage $t$}
    \State Recover $(S_t, Q_t)$ from $(t, \bar X_t^{(k_1)}, L^{(k_2)})$
    \State Compute mixture weights $\{\tilde w_k\}_{k=1}^M$ from $p(\sigma \mid X_{1:t})$
    \State Compute component means/variances $\{(\mu_t^{(\sigma_k)}, v_t^{(\sigma_k)})\}_{k=1}^M$ via~\eqref{eq:mu-static}
    \State $S \gets 0$
    \For{$k = 1, \ldots, M$, $j = 1, \ldots, J$}
      \State $x_{kj} \gets \mu_t^{(\sigma_k)} + \sqrt{2 v_t^{(\sigma_k)}}\, z_j$
      \State $S_{t+1} \gets S_t + x_{kj}$;\quad $Q_{t+1} \gets Q_t + x_{kj}^2$
      \State Translate to $(\bar X_{t+1}, L_{t+1})$, bilinearly interpolate $\widehat C_{t+1}^{\mathrm{bilin}}$ on stage-$(t{+}1)$'s $\bar X$ grid, clipping at grid boundary
      \State $S \gets S + \tilde w_k \cdot (w_j/\sqrt\pi) \cdot \max\!\big\{x_{kj}, \widehat C_{t+1}^{\mathrm{bilin}}\big\}$
    \EndFor
    \State $\widehat C_t(\bar X^{(k_1)}, L^{(k_2)}) \gets S$
  \EndFor
\EndFor
\State \Return $\{\widehat C_t\}$
\end{algorithmic}
\end{algorithm}

\paragraph{Convergence check.} Solve at $(K_1, K_2, J) = (256, 256, 64)$ and at reference $(512, 512, 128)$; report (i) max absolute disagreement on $\widehat C$ over the grid, and (ii) action-label disagreement on $10^4$ test sequences per regime. Fix the labeling resolution after the operational disagreement falls below $10^{-3}$.

\subsection{Labeling}
\label{sec:labeling}

For each training distribution we precompute the corresponding ADP threshold table $\{\widehat C_t\}$ once: Algorithm~\ref{alg:adp-static} for $\mathcal{D}_1, \mathcal{D}_2, \mathcal{D}_3$ (1D in $\mu_t$, linear interpolation), Algorithm~\ref{alg:adp-random} for $\mathcal{D}_{\mathrm{disc}}, \mathcal{D}_{\mathrm{logu}}$ (2D in $(\bar X_t, L_t)$, bilinear interpolation). Both tables expose the same query: given the belief state at time $t$, return $\widehat C_t$.

For each sequence $X_{1:n}$ from training distribution $\mathcal{D}$, at each $t \in \{1, \ldots, n-1\}$ we update the belief state and query the table for $\widehat C_t$, then emit
\[
y_t^{\mathrm{cv}} \;:=\; \widehat C_t,
\qquad
y_t^{\mathrm{act}} \;:=\; \mathbf{1}[X_t \ge \widehat C_t].
\]
At $t = n$, $y_n^{\mathrm{act}} = 1$ deterministically and $y_n^{\mathrm{cv}} = 0$.

\section{Portfolio of Baseline Algorithms}
\label{sec:baselines}

We report several decision rules: oracular baselines that assume the test-regime $\sigma$ is known, data-only baselines that estimate $\sigma$ from $X_{1:t}$ (directly comparable to the trained model), and scale-free baselines.

\subsection{Known-$\mu$ and known-$\sigma$ thresholds}
\label{sec:known-mu-sigma}

If both $\mu$ and $\sigma$ were known, the offers $X_1, \ldots, X_n$ would be uninformative about $\mu$ and the optimal continuation value $c_t := \mathbb{E}[\max\{X_{t+1}, c_{t+1}\}]$ would be a deterministic sequence with $c_{n-1} = \mu$. To evaluate the recursion, $\mathbb{E}[\max\{X, a\}] = \mu + \sigma\psi((a - \mu)/\sigma)$ for $X \sim \mathcal{N}(\mu, \sigma^2)$, where $\psi(z) := z\Phi(z) + \varphi(z)$. The standardized continuation value $\eta_t := (c_t - \mu)/\sigma$ depends on neither $\mu$ nor $\sigma$:
\begin{equation}
\label{eq:eta}
\eta_{n-1} = 0, \qquad \eta_t = \psi(\eta_{t+1}), \qquad t = n-2, \ldots, 1.
\end{equation}
The known-$\mu$, known-$\sigma$ optimal rule is $\mathbf{1}[X_t \ge \mu + \sigma\eta_t]$. Baselines below plug different estimators of $\mu$ and $\sigma$ into this rule.

\subsection{Known-$\sigma$ baselines}

Evaluated using the true test-regime $\sigma$.

\paragraph{Plug-in.} $\pi_t^{\mathrm{plug\text{-}in}}(X_{1:t}; \sigma) = \mathbf{1}[X_t \ge \bar X_t + \sigma\eta_t]$. Replaces $\mu$ by the sample mean $\bar X_t$.

\paragraph{Prior-only.} $\pi_t^{\mathrm{prior}}(X_t; \sigma) = \mathbf{1}[X_t \ge \mu_0 + \sigma\eta_t]$. Replaces $\mu$ by the prior mean.

\paragraph{Myopic.} $\pi_t^{\mathrm{myopic}}(X_{1:t}; \sigma) = \mathbf{1}[X_t \ge \mu_t]$, with $\mu_t$ the conjugate posterior mean at known $\sigma$.

\paragraph{Per-regime Bayes-optimal oracle $\pi^\star_{\mathcal{D}_i}$.} The Bayes-optimal policy from Algorithm~\ref{alg:adp-static} at $\sigma = \sigma_i$. The highest expected payoff achievable on $\mathcal{D}_i$ by any online policy.

\subsection{Data-only baselines}

These see only $X_{1:t}$ and estimate $\sigma$ from data --- directly comparable to the trained model.

\paragraph{MAP-$\sigma$ plug-in.} Compute the posterior $p(\sigma \mid X_{1:t})$ via~\eqref{eq:marginal-loglik}, take its mode --- the maximum-a-posteriori estimate
\[
\hat\sigma_t^{\mathrm{MAP}} \;:=\; \arg\max_\sigma p(\sigma \mid X_{1:t}),
\]
and apply
\[
\pi_t^{\mathrm{plug\text{-}in,\,MAP}}(X_{1:t}) \;=\; \mathbf{1}[X_t \ge \bar X_t + \hat\sigma_t^{\mathrm{MAP}}\, \eta_t].
\]

\paragraph{MLE-$\sigma$ plug-in.} The maximum-likelihood estimate of $\sigma$ is the within-sequence sample standard deviation
\[
\hat\sigma_t^{\mathrm{MLE}} \;:=\; \sqrt{\hat\sigma_t^2} \;=\; \sqrt{\frac{1}{t-1}\!\left(Q_t - t\bar X_t^2\right)}.
\]
Apply
\[
\pi_t^{\mathrm{plug\text{-}in,\,MLE}}(X_{1:t}) \;=\; \mathbf{1}[X_t \ge \bar X_t + \hat\sigma_t^{\mathrm{MLE}}\, \eta_t].
\]

\paragraph{Random-variance Bayes-optimal oracle $\pi^\star_{\mathcal{D}}$.} The Bayes-optimal policy from Algorithm~\ref{alg:adp-random} for $\mathcal{D} \in \{\mathcal{D}_{\mathrm{disc}}, \mathcal{D}_{\mathrm{logu}}\}$. Strict upper bound on every online baseline that does not know $\sigma$, under the prior of training distribution $\mathcal{D}$.

\subsection{Scale-free baselines}

\paragraph{Secretary.} Skip the first $r := \lfloor n/e \rfloor$ observations, set $M_r := \max_{s \le r} X_s$, accept the first $X_t$ ($t > r$) with $X_t > M_r$; if none, accept $X_n$.

\paragraph{Offline (hindsight).} $\pi^{\mathrm{offline}}(X_{1:n}) = \arg\max_t X_t$. Not implementable online; $\mathbb{E}[\max_t X_t]$ upper-bounds every online policy. By scale-and-shift, $\max_t X_t = \mu + \sigma Z_n$ where $Z_n$ is the maximum of $n$ iid $\mathcal{N}(0, 1)$ variables, and since $\mathbb{E}[\mu] = \mu_0 = 0$,
\[
\mathbb{E}[\max_t X_t;\, \mathcal{D}_i] \;=\; \sigma_i \cdot \mathbb{E}[Z_n], \qquad \mathbb{E}[Z_n] \;=\; \int_{-\infty}^{\infty} z \cdot n\Phi(z)^{n-1}\varphi(z)\, dz.
\]
The integral has no elementary closed form but is evaluated to machine precision by 1D quadrature.

\section{Model and Training}
\label{sec:training}

\subsection{Architecture}

GPT-2 decoder backbone with $L = 8$ layers, embedding dimension $d_{\mathrm{emb}} = 128$, $M = 4$ heads, causal attention mask, no dropout. Roughly $1.6$M parameters. The standard GPT-2 language-modeling head is removed; the model is not trained to predict $X_{t+1}$ from $X_{1:t}$. Each scalar $X_t$ is mapped to the embedding by a learnable affine layer $\mathbf{e}_t = W_{\mathrm{in}}\, X_t + \mathbf{b}_{\mathrm{in}}$, followed by learned positional embeddings. The hidden state $\mathbf{h}_t$ at each position feeds either a continuation-value head $\widehat C_t = \mathbf{w}_{\mathrm{cv}}^\top \mathbf{h}_t + b_{\mathrm{cv}}$ or an action head $\hat a_t = \sigma(\mathbf{w}_{\mathrm{act}}^\top \mathbf{h}_t + b_{\mathrm{act}})$, trained against the oracle labels (Section~\ref{sec:labeling}). The causal mask enforces that the output at position $t$ depends only on $X_{1:t}$, matching the online structure of the stopping problem. Inputs $X_t$ enter raw; the cv head produces $\widehat C_t$ in raw oracle units. No input-scale normalization is applied.

\subsection{Training objective}
\label{sec:training-objective}

We train ten models: five training distributions $\{\mathcal{D}_1, \mathcal{D}_2, \mathcal{D}_3, \mathcal{D}_{\mathrm{disc}}, \mathcal{D}_{\mathrm{logu}}\}$ $\times$ two supervisions $\{\text{cv}, \text{act}\}$. The model emits one prediction per timestep $t$ via the causal mask. We supervise at every $t \in \{1, \ldots, n-1\}$ (the terminal $t = n$ has no decision, since the agent must accept $X_n$): each timestep contributes its own oracle label and its own gradient signal, giving $n-1 = 255$ training signals per sequence. The per-sequence losses, with sequence-specific $\sigma_i$ used only for cv-loss normalization (the model never sees $\sigma_i$), are
\begin{align*}
\ell_{\mathrm{cv}}(\theta;\, X_{1:n},\, \sigma_i)
&\;=\; \frac{1}{(n-1)\,\sigma_i}\sum_{t=1}^{n-1}\!\Big(\widehat C_t(X_{1:t};\theta) - y_t^{\mathrm{cv}}\Big)^{\!2}, \\
\ell_{\mathrm{act}}(\theta;\, X_{1:n})
&\;=\; -\frac{1}{n-1}\sum_{t=1}^{n-1}\!\Big[\, y_t^{\mathrm{act}}\,\log \hat a_t + (1 - y_t^{\mathrm{act}})\,\log(1 - \hat a_t)\,\Big].
\end{align*}
Each is the within-sequence mean over per-step terms. The mini-batch loss is the further mean over $B = 64$ sequences sampled fresh from the training distribution.

Without normalization, cv squared errors have typical magnitude $\sim \sigma^2$, so a $\sigma = 100$ batch contributes gradients $10^4 \times$ larger than a $\sigma = 1$ batch and a single Adam learning rate cannot accommodate both. We divide by $\sigma_i$ rather than $\sigma_i^2$ to make per-sequence gradient magnitudes regime-invariant: at any step the residual $\widehat C_t - y_t^{\mathrm{cv}}$ has typical scale $\sigma_i$, so $d\ell_{\mathrm{cv}}/d\theta \propto (1/\sigma_i)\cdot\text{residual}\cdot d\widehat C_t/d\theta$ is order unity in $\sigma_i$.\footnote{$1/\sigma_i$ instead of $1/\sigma_i^2$ gives gradient invariance across $\sigma$ regimes; loss values now scale as $\sigma$ and are not directly comparable across regimes, but loss-value comparability across regimes was already not held under $1/\sigma_i^2$ either (only each model's own descent shape is meaningful as a convergence proxy). Logged in the V3 README's spec corrections.} The normalizer is per-sequence (constant within a sequence) and is applied on the loss side only: it gives the model no in-sequence information about $\sigma_i$ --- the model still has to recover $\sigma$ from $X_{1:t}$ on its own. Act supervision is naturally regime-invariant. One consequence: cv val MSE values are in $\sigma^{-1}$ units and not comparable across regimes; only each model's own MSE descent indicates convergence.

Adam, peak learning rate $10^{-4}$, no weight decay, batch size $64$. Static-variance cv models train for $2 \times 10^5$ steps; random-variance cv models for $3 \times 10^5$ (additional steps for in-context regime estimation). Act models train for half their cv counterpart's step count. Cosine schedule with $20\%$ linear warmup; learned absolute position embeddings; checkpoint selection by minimum validation loss.

\section{In-Context Learning: What and How?}
\label{sec:experiments}

All ten models are evaluated on each of the three test regimes $\mathcal{D}_1, \mathcal{D}_2, \mathcal{D}_3$. We split the analysis into two experiments: the six static-variance models (Section~\ref{sec:exp-static}) and the four random-variance models (Section~\ref{sec:exp-random}).

\subsection{Metrics}

For each (model, test regime) pair we report expected payoff $R(\pi; \mathcal{D}_i) := \mathbb{E}[X_{\tau_\pi}]$, normalized payoff $R/R^\star_{\mathcal{D}_i}$ relative to the per-regime oracle (the regime-invariant metric for cross-regime comparison), and per-step action agreement with each baseline. We report agreement as a diagnostic --- two policies that reject most positions can agree on $> 98\%$ of steps while differing on the few accepts that determine the entire payoff, so high agreement does not imply matched payoff --- but $R/R^\star$ is the headline metric.

\subsection{Static-variance experiment}
\label{sec:exp-static}

The six static-variance models test whether single-regime training induces a regime-invariant Bayesian algorithm or merely a fixed-scale policy specific to the training regime. Each model is trained on one of $\mathcal{D}_1, \mathcal{D}_2, \mathcal{D}_3$ and evaluated on all three test regimes, with all baselines reported per (model, test-regime) cell. The off-diagonal $R/R^\star$ values are the key diagnostic.

\subsection{Random-variance experiment}
\label{sec:exp-random}

The four random-variance models test whether training on a non-degenerate $\sigma$-prior induces in-context regime inference. Each model is trained on one of $\mathcal{D}_{\mathrm{disc}}, \mathcal{D}_{\mathrm{logu}}$ and evaluated on all three test regimes, with all baselines reported per cell.

\paragraph{Comparison to data-only baselines.} The MAP-$\sigma$ and MLE-$\sigma$ plug-in baselines see exactly the same input as the model: $X_{1:t}$. The MAP-$\sigma$ baseline uses the same prior over $\sigma$ as the model's training distribution. The trained model can outperform these baselines either by being closer to its training-distribution oracle $\pi^\star_{\mathcal{D}}$, or by exploiting features of $X_{1:t}$ beyond $(\bar X_t, \hat\sigma_t^2)$.

\subsection{Threshold trajectories}
\label{sec:exp-thresh}

For each cv-supervised model and each test regime, plot the per-step threshold (mean $\pm 1$ std over the $N_{\mathrm{test}}$ test sequences) against the per-regime oracle, the random-variance oracle (where applicable), the data-only baselines (MAP-$\sigma$, MLE-$\sigma$), the known-$\sigma$ baselines, the secretary baseline (the running max $M_r$ over the first $r = \lfloor n/e \rfloor$ observations, treated as $-\infty$ for $t \le r$ and held constant for $t > r$), and the offline-hindsight $\mathbb{E}[\max_{1 \le s \le n} X_s; \mathcal{D}_i]$ as a horizontal line per regime (computed via Section~\ref{sec:baselines}). A model that has learned in-context $\sigma$ inference produces threshold trajectories that adapt to the test regime within a few steps; one that has not produces a flat regime-invariant threshold.

\section{Results}
\label{sec:results}

\subsection{Static-variance models reproduce the V2 shortcut}
\label{sec:res-static-shortcut}

Each single-regime model matches the per-regime oracle in-distribution and degrades sharply off-diagonally. Table~\ref{tab:res-static} reports $R/R^\star$ for the three static-variance cv runs on the three test regimes; the off-diagonal cells fall to as low as $0.016$ ($\mathcal{D}_3$\_cv on $\mathcal{D}_1$).

\begin{table}[H]
\centering
\small
\begin{tabular}{lccc}
\hline
trained model & $\mathcal{D}_1$ & $\mathcal{D}_2$ & $\mathcal{D}_3$ \\
\hline
$\mathcal{D}_1$\_cv & $1.000$ & $0.318$ & $0.309$ \\
$\mathcal{D}_2$\_cv & $0.174$ & $1.000$ & $0.362$ \\
$\mathcal{D}_3$\_cv & $0.016$ & $0.052$ & $1.000$ \\
\hline
\end{tabular}
\caption{$R/R^\star$ for the static-variance cv-supervised models across the three test regimes. Diagonal entries are in-distribution; off-diagonal entries are out-of-distribution. The act counterparts are within $0.04$ of the cv values cell-for-cell.}
\label{tab:res-static}
\end{table}

This is the V2 single-regime shortcut, reproduced under the V3 architecture. V2 employed input-side scale buffers to handle multi-regime $\sigma$; V3 has no such buffers (Section~\ref{sec:training}: $X_t$ enters raw, $\widehat C_t$ exits raw), and the static models still collapse off-diagonally. The shortcut is therefore a property of single-regime training rather than an artifact of architectural normalization choice.

Figures~\ref{fig:traj-d1cv}--\ref{fig:traj-d3cv} visualize the same effect at the threshold-trajectory level: each model's mean threshold sits flat at its training-distribution scale across all three test regimes, with no adaptation observable in the test prefix. This is the trajectory-level signature of a memorized fixed-scale policy.

\begin{figure}[H]
\centering
\includegraphics[width=\textwidth]{figures/trajectories-d-1-cv.png}
\caption{Threshold trajectories for $\mathcal{D}_1$\_cv (focal model, red, mean $\pm 1$ std band over $N_{\mathrm{test}} = 10^4$ sequences) on each of the three test regimes, with all baselines from Section~\ref{sec:baselines} overlaid: the per-regime oracle, the random-variance oracles for $\mathcal{D}_{\mathrm{disc}}$ and $\mathcal{D}_{\mathrm{logu}}$, the known-$\sigma$ rules (plug-in, prior-only, myopic), the data-only plug-ins (MAP-$\sigma$ at both priors, MLE-$\sigma$), the secretary running max, and the offline-hindsight expected payoff (horizontal line). The model's threshold remains at the $\sigma = 1$ scale ($\approx 2.5$) on every regime.}
\label{fig:traj-d1cv}
\end{figure}

\begin{figure}[H]
\centering
\includegraphics[width=\textwidth]{figures/trajectories-d-2-cv.png}
\caption{Threshold trajectories for $\mathcal{D}_2$\_cv. The model's threshold remains at the $\sigma = 10$ scale ($\approx 26$) on every regime.}
\label{fig:traj-d2cv}
\end{figure}

\begin{figure}[H]
\centering
\includegraphics[width=\textwidth]{figures/trajectories-d-3-cv.png}
\caption{Threshold trajectories for $\mathcal{D}_3$\_cv. The model's threshold remains at the $\sigma = 100$ scale ($\approx 260$) on every regime.}
\label{fig:traj-d3cv}
\end{figure}

Figure~\ref{fig:traj-all-cv} overlays all five cv-trained models on each test regime, juxtaposing the three static-variance plateaus against the multi-regime $\mathcal{D}_{\mathrm{disc}}$\_cv adaptation; the contrast is visible at a glance, and previews the $\mathcal{D}_{\mathrm{disc}}$ result analyzed in Section~\ref{sec:res-disc}.

\begin{figure}[H]
\centering
\includegraphics[width=\textwidth]{figures/trajectories-all-cv.png}
\caption{All five cv-trained models overlaid on each test regime, with the per-regime oracle (black dashed) for reference. The three single-regime models ($\mathcal{D}_1$\_cv red, $\mathcal{D}_2$\_cv orange, $\mathcal{D}_3$\_cv purple) plateau at their training-distribution scale on every regime, regardless of the data they observe; $\mathcal{D}_{\mathrm{disc}}$\_cv (blue, thicker) tracks the per-regime oracle on each panel; $\mathcal{D}_{\mathrm{logu}}$\_cv (green) tracks on $\mathcal{D}_3$ but undershoots elsewhere (Section~\ref{sec:res-logu-cv}). Multi-regime training is what enables in-context regime adaptation; single-regime training does not.}
\label{fig:traj-all-cv}
\end{figure}

\subsection{$\mathcal{D}_{\mathrm{disc}}$ training recovers the regime-invariant algorithm}
\label{sec:res-disc}

Multi-regime training under $\mathcal{D}_{\mathrm{disc}}$ removes the single-regime shortcut. Both $\mathcal{D}_{\mathrm{disc}}$\_cv and $\mathcal{D}_{\mathrm{disc}}$\_act match the per-regime oracle on every test regime: $R/R^\star = 0.999, 0.999, 1.000$ for cv and $0.999, 0.999, 0.999$ for act, on $\mathcal{D}_1, \mathcal{D}_2, \mathcal{D}_3$ respectively (Figure~\ref{fig:payoff-matrix}). A single trained model attains, on every regime, what each single-regime specialist attains only in-distribution --- the shortcut has been replaced by a regime-invariant policy.

\begin{figure}[H]
\centering
\includegraphics[width=0.55\textwidth]{figures/payoff-matrix.png}
\caption{Normalized expected payoff $R/R^\star$ for every (trained model, test regime) cell. Rows are the ten trained models; columns are the three test regimes ($\mathcal{D}_1$, $\mathcal{D}_2$, $\mathcal{D}_3$). $R^\star$ is the per-regime static-$\sigma$ oracle (Algorithm~\ref{alg:adp-static}).}
\label{fig:payoff-matrix}
\end{figure}

The per-$\sigma$ payoff breakdown on the seed-$44$ $\mathcal{D}_{\mathrm{disc}}$ test cache (Figure~\ref{fig:per-sigma-disc}) confirms uniform competence across the training prior: $R/R^\star \in \{1.001, 0.999, 0.999\}$ for $\mathcal{D}_{\mathrm{disc}}$\_cv at $\sigma \in \{1, 10, 100\}$, and within $0.001$ of the oracle for $\mathcal{D}_{\mathrm{disc}}$\_act in every bin. The per-$\sigma$ test cache draws sequences from the training distribution and reports $R/R^\star$ against the per-distribution random-ADP oracle (Algorithm~\ref{alg:adp-random}) bin-wise; restricting to a $\sigma$-bin isolates per-$\sigma$ competence from $\sigma$-mixing artifacts that aggregated $R/R^\star$ on the regime test caches would absorb.

\begin{figure}[H]
\centering
\includegraphics[width=0.85\textwidth]{figures/per-sigma-d-disc.png}
\caption{Per-$\sigma$ payoff breakdown for $\mathcal{D}_{\mathrm{disc}}$\_cv (blue) and $\mathcal{D}_{\mathrm{disc}}$\_act (orange) on the seed-$44$ per-$\sigma$ test cache. Both models track the per-distribution random-ADP oracle ($R/R^\star = 1$, dashed) in every $\sigma$ bin.}
\label{fig:per-sigma-disc}
\end{figure}

The threshold trajectory for $\mathcal{D}_{\mathrm{disc}}$\_cv (Figure~\ref{fig:traj-disc-cv}) shows the corresponding signature: on each test regime, the model's threshold moves toward the per-regime oracle within a few observations of $X_t$ and tracks it for the remainder of the horizon. The model is performing in-context $\sigma$ inference and applying the regime-corresponding threshold rule --- the algorithm hypothesis (Section~\ref{sec:intro}) holds when the training distribution provides examples spanning multiple $\sigma$ regimes.

\begin{figure}[H]
\centering
\includegraphics[width=\textwidth]{figures/trajectories-d-disc-cv.png}
\caption{Threshold trajectory for $\mathcal{D}_{\mathrm{disc}}$\_cv (focal, blue, mean $\pm 1$ std band) on each test regime, with all baselines from Section~\ref{sec:baselines} overlaid. The model's threshold adapts to the test regime within a few observations and tracks the per-regime oracle (black dashed) and the $\mathcal{D}_{\mathrm{disc}}$ random-ADP oracle (its in-distribution Bayes-optimal reference) for the remainder of the horizon on every regime.}
\label{fig:traj-disc-cv}
\end{figure}

\subsection{$\mathcal{D}_{\mathrm{logu}}$\_cv exhibits a supervision-form-dependent failure}
\label{sec:res-logu-cv}

Training under the continuous wide prior $\mathcal{D}_{\mathrm{logu}}$ produces a markedly asymmetric outcome between the two supervision forms. $\mathcal{D}_{\mathrm{logu}}$\_cv attains $R/R^\star = 0.036, 0.091, 0.800$ on $\mathcal{D}_1, \mathcal{D}_2, \mathcal{D}_3$ --- monotonically degrading toward small $\sigma$ --- while $\mathcal{D}_{\mathrm{logu}}$\_act recovers fully, attaining $R/R^\star = 0.974, 0.976, 0.975$ on the same three regimes (Figure~\ref{fig:payoff-matrix}). Figure~\ref{fig:traj-logu-cv} shows the corresponding trajectory: $\mathcal{D}_{\mathrm{logu}}$\_cv's mean threshold sits well above the per-regime oracle on $\mathcal{D}_1$ and $\mathcal{D}_2$ for the entire horizon and only approaches the oracle on $\mathcal{D}_3$. Figure~\ref{fig:traj-zoom} re-renders the same comparison restricted to $\mathcal{D}_{\mathrm{disc}}$\_cv and $\mathcal{D}_{\mathrm{logu}}$\_cv at per-panel y-limits chosen to expose the residual gap on each regime.

\begin{figure}[H]
\centering
\includegraphics[width=\textwidth]{figures/trajectories-d-logu-cv.png}
\caption{Threshold trajectory for $\mathcal{D}_{\mathrm{logu}}$\_cv (focal, green, mean $\pm 1$ std band) on each test regime, with all baselines from Section~\ref{sec:baselines} overlaid.}
\label{fig:traj-logu-cv}
\end{figure}

\begin{figure}[H]
\centering
\includegraphics[width=\textwidth]{figures/trajectories-disc-logu-zoom.png}
\caption{Zoomed comparison of $\mathcal{D}_{\mathrm{disc}}$\_cv (blue) and $\mathcal{D}_{\mathrm{logu}}$\_cv (green) against the per-regime oracle (black dashed) on each test regime, with per-panel y-limits chosen to make the residual gap visible: $y \in [0, 10]$ on $\mathcal{D}_1$, $[0, 40]$ on $\mathcal{D}_2$, $[150, 300]$ on $\mathcal{D}_3$.}
\label{fig:traj-zoom}
\end{figure}

The per-$\sigma$ payoff breakdown on the $\mathcal{D}_{\mathrm{logu}}$ test cache (Figure~\ref{fig:per-sigma-logu}) confirms that the regime-level effect is monotone in $\sigma$ rather than a coincidence of where the three test regimes land: $R/R^\star$ for $\mathcal{D}_{\mathrm{logu}}$\_cv climbs from $-0.044$ in the smallest $\log_{10}\sigma$ bin to $0.722$ in the largest, while $\mathcal{D}_{\mathrm{logu}}$\_act stays within $[0.998, 1.000]$ across all five bins.

\begin{figure}[H]
\centering
\includegraphics[width=0.95\textwidth]{figures/per-sigma-d-logu.png}
\caption{Per-$\sigma$ payoff breakdown for $\mathcal{D}_{\mathrm{logu}}$\_cv (blue) and $\mathcal{D}_{\mathrm{logu}}$\_act (orange) on the seed-$44$ per-$\sigma$ test cache. The cv model degrades monotonically toward small $\sigma$; the act model is uniformly competitive with the per-distribution random-ADP oracle.}
\label{fig:per-sigma-logu}
\end{figure}

The supervision form is the central observation. Both models share the architecture (Section~\ref{sec:training}), the training distribution, the data, and the training step count; only the per-step target differs --- the continuation value $\widehat C_t$ for cv, the binary action $\mathbf{1}[X_t \ge \widehat C_t]$ for act. Action supervision recovers fully; continuation-value supervision does not.

To localize the failure within the cv model we evaluate, on the same per-$\sigma$ test cache, the discrepancy between the model's threshold output $\widehat C_t$ and the bin's static-$\sigma$ oracle threshold $C_t^\star$ at the bin's geometric-mean $\sigma$, averaged over the cv-mask range $t \in \{2, \ldots, n - 1\}$. Comparing in two units --- raw threshold (Figure~\ref{fig:diag-abs}, left) and the loss's natural units $|\widehat C_t - C_t^\star|/\sigma_i$ (Figure~\ref{fig:diag-abs}, right) --- yields opposite scaling.

\begin{figure}[H]
\centering
\begin{minipage}{0.49\textwidth}
\centering
\includegraphics[width=\textwidth]{figures/diagnostic-logu-abs-error.png}
\end{minipage}\hfill
\begin{minipage}{0.49\textwidth}
\centering
\includegraphics[width=\textwidth]{figures/diagnostic-logu-loss-unit-error.png}
\end{minipage}
\caption{Decomposition of the $\mathcal{D}_{\mathrm{logu}}$\_cv error by $\log_{10}\sigma$ bin, on the seed-$44$ per-$\sigma$ test cache. Left: mean absolute error $|\widehat C_t - C_t^\star|$ in raw threshold units. Right: the same error in the cv loss's natural units, $|\widehat C_t - C_t^\star|/\sigma_i$.}
\label{fig:diag-abs}
\end{figure}

In raw threshold units, $|\widehat C_t - C_t^\star|$ grows from $0.871$ in the smallest $\sigma$ bin to $32.296$ in the largest --- a $\sim 37\times$ range, roughly proportional to $\sigma$ across the prior. In loss units the same quantity is essentially flat: $0.539, 0.562, 0.544, 0.547, 0.525$ across the five bins, a $1.07\times$ ratio between maximum and minimum. Read together, the two scalings show that the cv loss normalization (Section~\ref{sec:training-objective}, $\ell_{\mathrm{cv}} \propto 1/\sigma_i$) is functioning as designed: the trained model achieves $\sigma$-invariant residual error in the units the loss measures, the gradient invariance the normalizer was intended to deliver. The model's internal representation is uniformly competent across the training prior.

The bottleneck is therefore in the cv head, not the representation. The cv head is a single linear map $\mathbf{w}_{\mathrm{cv}}^\top \mathbf{h}_t + b_{\mathrm{cv}}$. A constant $\sigma$-relative error in $\widehat C_t$ corresponds to a constant residual in loss space but to a $\sigma$-proportional error in raw threshold units, and downstream to a stopping-rule error whose payoff impact depends on the absolute scale of $C_t^\star$. At $\sigma = 1$ the optimal threshold is $\sim 2.5$ and the raw error is $\sim 0.87$, of comparable magnitude to the threshold itself, so the comparison $X_t \ge \widehat C_t$ flips on a substantial fraction of decision-critical timesteps. At $\sigma = 100$ the optimal threshold is $\sim 260$ and the raw error is $\sim 32$, sufficiently smaller than the threshold that the comparison's outcome is mostly preserved. The loss-balanced training therefore converts into a downstream payoff that scales sharply with $\sigma$.

Action supervision is robust to this failure mode because the per-step target is the binary label $\mathbf{1}[X_t \ge \widehat C_t]$ rather than $\widehat C_t$ itself. The act head emits a logit, BCE penalizes wrong actions, and the gradient signal directly reinforces decision-correct behavior. Provided the model's internal threshold is in the right ballpark for the test regime, the comparison $X_t \ge \widehat C_t$ resolves correctly even when $\widehat C_t$ as a number would carry residual noise that, exposed through the linear cv head, would translate into payoff loss. The failure mode of cv supervision is invisible to act supervision because act supervision never reads $\widehat C_t$ as a number.

This is therefore a structural property of regression-through-a-linear-head supervision over a wide-prior task, not an engineering shortfall of the V3 setup. The cv loss normalization works as designed; the failure is downstream of the loss. Mitigations are conceivable --- a non-linear continuation-value head, a log-target regression formulation, or simply act supervision when the prior on $\sigma$ spans multiple decades --- but their evaluation is deferred to future work.

\subsection{Comparison to data-only baselines}
\label{sec:res-baselines}

The MAP-$\sigma$ and MLE-$\sigma$ plug-in baselines (Section~\ref{sec:baselines}) see exactly the same input as the trained random-variance models. Comparing the trained models against these baselines isolates whether the model is adding value beyond data-only $\sigma$-estimation followed by formula plug-in.

\begin{table}[H]
\centering
\small
\begin{tabular}{lccc}
\hline
policy & $\mathcal{D}_1$ & $\mathcal{D}_2$ & $\mathcal{D}_3$ \\
\hline
$\mathcal{D}_{\mathrm{disc}}$\_cv                         & $0.999$ & $0.999$ & $1.000$ \\
$\mathcal{D}_{\mathrm{disc}}$\_act                        & $0.999$ & $0.999$ & $0.999$ \\
$\mathcal{D}_{\mathrm{logu}}$\_cv                         & $0.036$ & $0.091$ & $0.800$ \\
$\mathcal{D}_{\mathrm{logu}}$\_act                        & $0.974$ & $0.976$ & $0.975$ \\
MAP-$\sigma$ ($\mathcal{D}_{\mathrm{disc}}$ prior)        & $1.000$ & $0.999$ & $0.991$ \\
MAP-$\sigma$ ($\mathcal{D}_{\mathrm{logu}}$ prior)        & $0.974$ & $0.973$ & $0.975$ \\
MLE-$\sigma$ plug-in                                      & $0.978$ & $0.976$ & $0.970$ \\
\hline
\end{tabular}
\caption{$R/R^\star$ of random-variance trained models versus data-only plug-in baselines, evaluated on the seed-$43$ test caches.}
\label{tab:res-baselines}
\end{table}

$\mathcal{D}_{\mathrm{disc}}$\_cv and $\mathcal{D}_{\mathrm{disc}}$\_act outperform every data-only baseline on every regime --- in particular the MAP-$\sigma$ baseline that uses the same prior over $\sigma$ as the model's training distribution. The margin is small but consistent ($1.000$ versus $0.991$ on $\mathcal{D}_3$, where the prior on $\sigma$ is most informative), evidence that the trained $\mathcal{D}_{\mathrm{disc}}$ model is computing a feature of $X_{1:t}$ beyond $(\bar X_t, \hat\sigma_t^2)$ that the closed-form plug-in does not exploit. $\mathcal{D}_{\mathrm{logu}}$\_cv, in contrast, loses to both data-only baselines on $\mathcal{D}_1$ ($0.036$ versus $\ge 0.974$) and on $\mathcal{D}_2$ ($0.091$ versus $\ge 0.973$). This is consistent with the head-bottleneck localization in Section~\ref{sec:res-logu-cv}: the data-only baselines emit thresholds via closed-form formulae that do not pass through the trained model's linear cv head, so they are not subject to the same head-induced threshold noise. $\mathcal{D}_{\mathrm{logu}}$\_act recovers to within $0.001$ of the MAP-$\sigma$ ($\mathcal{D}_{\mathrm{logu}}$ prior) baseline on every regime, consistent with action supervision's robustness to that noise.

\subsection{Agreement-vs-payoff pathology}
\label{sec:res-agreement}

Per-step action agreement and normalized payoff can differ sharply on the same (model, test regime) cell. $\mathcal{D}_3$\_cv on $\mathcal{D}_1$ agrees with the per-regime oracle on $97.8\%$ of its decisions yet attains $R/R^\star = 0.016$; $\mathcal{D}_{\mathrm{logu}}$\_cv on $\mathcal{D}_1$ agrees on $99.5\%$ of decisions yet attains $R/R^\star = 0.036$. Figure~\ref{fig:agreement-oracle} shows the full ten-by-three matrix. The pathology is mechanical: most timesteps in a sequence are trivial rejects --- the optimal action is ``continue'' and any policy emitting a sufficiently large threshold also rejects --- which inflates the agreement metric. The few decision-critical timesteps where the optimal action is ``accept'' determine the entire payoff, and a fixed-scale threshold flips precisely those decisions in the wrong direction. The phenomenon was identified in V2's analysis of the static-variance shortcut and reproduces here under the V3 architecture. We retain agreement as a diagnostic accompanying $R/R^\star$, never as a substitute for it.

\begin{figure}[H]
\centering
\includegraphics[width=0.5\textwidth]{figures/agreement-oracle.png}
\caption{Per-step action agreement of each trained model with the per-regime oracle on each test regime. High agreement is compatible with very low $R/R^\star$ (e.g., $\mathcal{D}_3$\_cv on $\mathcal{D}_1$: $0.978$ agreement, $0.016$ payoff; $\mathcal{D}_{\mathrm{logu}}$\_cv on $\mathcal{D}_1$: $0.995$ agreement, $0.036$ payoff).}
\label{fig:agreement-oracle}
\end{figure}

\subsection{Findings}
\label{sec:res-findings}

\begin{itemize}
\item Multi-regime training under the discrete $\sigma$-prior $\mathcal{D}_{\mathrm{disc}}$ recovers the regime-invariant Bayesian algorithm: $\mathcal{D}_{\mathrm{disc}}$\_cv and $\mathcal{D}_{\mathrm{disc}}$\_act both attain $R/R^\star \ge 0.999$ on every test regime, and threshold adaptation occurs within a few observations of $X_t$.

\item The continuous $\sigma$-prior $\mathcal{D}_{\mathrm{logu}}$ reveals a supervision-form-dependent failure. Continuation-value supervision degrades monotonically toward small $\sigma$ ($R/R^\star = 0.036, 0.091, 0.800$ on $\mathcal{D}_1, \mathcal{D}_2, \mathcal{D}_3$); action supervision recovers fully ($R/R^\star \approx 0.97$ on every regime).

\item An error decomposition localizes the $\mathcal{D}_{\mathrm{logu}}$\_cv failure to the linear continuation-value head, not the model's representation. Loss-unit error $|\widehat C_t - C_t^\star|/\sigma_i$ is constant across $\sigma$ ($0.54 \pm 0.02$, $1.07\times$ ratio between maximum and minimum bin), confirming that the cv loss normalization functions as designed. The failure is structural, not a tuning artifact.

\item Trained random-variance models on $\mathcal{D}_{\mathrm{disc}}$ outperform data-only plug-in baselines on every regime, evidence the model is doing more than $\sigma$-estimation followed by formula plug-in. $\mathcal{D}_{\mathrm{logu}}$\_cv loses to data-only baselines on $\mathcal{D}_1$ and $\mathcal{D}_2$, consistent with the head-bottleneck localization: the data-only baselines emit thresholds via formulae that do not pass through the trained linear head.

\item V2's static-variance shortcut reproduces faithfully under the V3 architecture, which uses no input-side scale buffers. The shortcut is a property of single-regime training, not an artifact of architectural normalization choice.
\end{itemize}

\section{Future work --- interpretability suite}
\label{sec:future}

The interpretability analyses below were planned alongside the main experiments and are deferred to follow-up work. Each is scoped against a specific question raised by Section~\ref{sec:results} --- which makes the deferred work narrower and more concrete than it would have been before the §6 results.

\subsection{Attention entropy}
\label{sec:fw-attn}

Normalized Shannon-entropy summary of attention distributions across the test set, plotted by (layer, head), to identify candidate running-mean and non-trivial-computation heads. Deferred: Section~\ref{sec:res-logu-cv} localizes the $\mathcal{D}_{\mathrm{logu}}$\_cv failure to the linear cv head; an interpretability follow-up would ask whether the same head pattern carries the noise level in $\mathcal{D}_{\mathrm{disc}}$ models, and where in the residual stream $\sigma$ becomes decodable.

\subsection{Probing}
\label{sec:fw-probe}

Auxiliary probe networks (linear and 2-layer MLP) trained to recover sufficient statistics ($S_t$, $\bar X_t$, $Q_t$, $\hat\sigma_t^2$), point estimates of $\sigma$ ($\hat\sigma_t^{\mathrm{MAP}}$, $\hat\sigma_t^{\mathrm{MLE}}$), and the oracle threshold $\widehat C_t$ from frozen-model hidden states at each layer; reported as probe MSE plus a learned position-attention readout. Deferred: a probing follow-up would test whether the uniformly-competent representation diagnosed in Section~\ref{sec:res-logu-cv} actually carries a recoverable $\sigma$ estimate at the layer immediately before the cv head.

\subsection{Causal interventions}
\label{sec:fw-intervention}

Vector arithmetic on the residual stream to test causal use of probe-direction-encoded quantities: push the layer-$\ell$ residual along a probe direction $\mathbf{w}$ and report the resulting shift in $\widehat C_t$ on a held-out test set, normalized by the natural between-regime shift. Deferred: causal use of $\hat\sigma_t^2$ or $\hat\sigma_t^{\mathrm{MAP}}$ becomes interpretable only after Section~\ref{sec:results} fixes a specific failure (the cv head) and a specific control (act supervision) to compare against.

\subsection{Logit lens for the continuation-value head}
\label{sec:fw-logit-lens}

Apply the trained cv head (or act head) to intermediate-layer hidden states to read out $\widehat C_t^{(\ell)}$ at every depth, identifying the layer at which the model commits to its threshold. Deferred: combined with probing, this would order the depths at which sufficient statistics, $\sigma$ estimates, and the threshold itself become decodable --- structurally informative once Section~\ref{sec:res-logu-cv}'s diagnosis identifies which depth's representation is being mis-used by the head.

\subsection{Activation patching}
\label{sec:fw-patching}

Replace the activation at a candidate site (layer $\ell$, position $t'$, component class $c \in \{\text{attention}, \text{MLP}, \text{head}\;m\}$) on a $\mathcal{D}_2$ ``clean'' sequence with the corresponding activation from a $\mathcal{D}_3$ ``corrupt'' sequence, and measure the resulting shift in $\widehat C_t$. Deferred: activation patching localizes the §6 phenomena to specific sites; the §6 results pin which phenomena to localize first.

\subsection{Head ablation at inference time}
\label{sec:fw-ablation}

Zero or mean-replace one attention head's output at inference and report the resulting $R/R^\star$ on each test regime; per-head, per-regime drop $\Delta_{\ell,m}^{(i)}$ identifies regime-specialized vs.\ regime-invariant heads. Deferred: Section~\ref{sec:results} identifies which models harbor regime-specialized circuitry worth ablating ($\mathcal{D}_{\mathrm{disc}}$\_cv, $\mathcal{D}_{\mathrm{logu}}$\_act); cheap to run once those targets are fixed.

\appendix

\section{Training curves}
\label{app:training}

Training ran for $2 \times 10^5$ cv steps and $1 \times 10^5$ act steps on the static-variance models, and $3 \times 10^5$ cv steps and $1.5 \times 10^5$ act steps on the random-variance models, with the additional budget on the random-variance side allotted to in-context regime estimation (Section~\ref{sec:training-objective}). Adam, peak learning rate $10^{-4}$, no weight decay, batch size $64$, cosine schedule with $20\%$ linear warmup. Figure~\ref{fig:training-curves} shows train and validation loss for all ten runs over their training horizon, with the final-checkpoint step (selected by minimum validation loss) marked.

\begin{figure}[H]
\centering
\includegraphics[width=\textwidth]{figures/training-curves.png}
\caption{Train (light) and validation (dark) loss for all ten trained runs, with the best-validation-loss checkpoint step marked (red dashed).}
\label{fig:training-curves}
\end{figure}

For the four random-variance runs we additionally report validation loss separated by $\sigma$-bin (Figure~\ref{fig:training-curves-per-sigma}), the diagnostic for the cv-loss normalization choice in Section~\ref{sec:training-objective}. On $\mathcal{D}_{\mathrm{disc}}$\_cv the $\sigma = 100$ residual lags the $\sigma = 1$ residual through most of training and converges to within $\sim 0.05$ of the $\sigma = 1$ residual by the end of the $3 \times 10^5$-step budget --- the empirical basis for the $1/\sigma_i$ normalization (rather than $1/\sigma_i^2$) used in $\ell_{\mathrm{cv}}$. The same catchup pattern is visible on $\mathcal{D}_{\mathrm{logu}}$\_cv.

\begin{figure}[H]
\centering
\includegraphics[width=\textwidth]{figures/training-curves-per-sigma.png}
\caption{Per-$\sigma$-bin validation loss for the four random-variance training runs ($\mathcal{D}_{\mathrm{disc}}$\_cv, $\mathcal{D}_{\mathrm{disc}}$\_act, $\mathcal{D}_{\mathrm{logu}}$\_cv, $\mathcal{D}_{\mathrm{logu}}$\_act). Black: aggregate; coloured: per-$\sigma$ slice.}
\label{fig:training-curves-per-sigma}
\end{figure}

\section{Derivation of the marginal log-likelihood}
\label{app:marginal-derivation}

We derive equation~\eqref{eq:marginal-loglik} by simplifying the determinant and inverse of the covariance $\Sigma_t^{(\sigma)} = \sigma^2 I_t + \tau_0^2 \mathbf{1}\mathbf{1}^\top$. Since $\Sigma_t^{(\sigma)}$ is a rank-one update of a scaled identity, both the matrix determinant lemma and the Sherman--Morrison formula apply: for invertible $A$ and column vectors $u, v$,
\[
\det(A + uv^\top) \;=\; \det(A) \cdot (1 + v^\top A^{-1} u),
\qquad
(A + uv^\top)^{-1} \;=\; A^{-1} - \frac{A^{-1} u v^\top A^{-1}}{1 + v^\top A^{-1} u}.
\]
In our case $A = \sigma^2 I_t$, $u = v = \tau_0 \mathbf{1}$, so $u v^\top = \tau_0^2 \mathbf{1}\mathbf{1}^\top$ and $\Sigma_t^{(\sigma)} = A + uv^\top$. The starting point is the log density of $X_{1:t} \mid \sigma \sim \mathcal{N}(\mathbf{0}, \Sigma_t^{(\sigma)})$,
\[
\log p(X_{1:t} \mid \sigma) \;=\; -\frac{t}{2}\log(2\pi) - \frac{1}{2}\log\det\Sigma_t^{(\sigma)} - \frac{1}{2} X_{1:t}^\top\!\big(\Sigma_t^{(\sigma)}\big)^{\!-1} X_{1:t}.
\]

For the determinant: $A^{-1} = \sigma^{-2} I_t$ gives $v^\top A^{-1} u = \tau_0^2 \mathbf{1}^\top \mathbf{1}/\sigma^2 = t\tau_0^2/\sigma^2$, and $\det A = \sigma^{2t}$, so by the matrix determinant lemma
\[
\log\det\Sigma_t^{(\sigma)} \;=\; t\log\sigma^2 + \log\!\left(1 + \frac{t\tau_0^2}{\sigma^2}\right).
\]
For the inverse: $A^{-1} u = \sigma^{-2}\tau_0\mathbf{1}$ and $A^{-1} u v^\top A^{-1} = \tau_0^2 \mathbf{1}\mathbf{1}^\top/\sigma^4$, so by Sherman--Morrison
\[
\big(\Sigma_t^{(\sigma)}\big)^{\!-1} \;=\; \frac{1}{\sigma^2}I_t - \frac{\tau_0^2/\sigma^4}{1 + t\tau_0^2/\sigma^2}\mathbf{1}\mathbf{1}^\top.
\]
The quadratic form follows from $X_{1:t}^\top I_t X_{1:t} = Q_t$ and $X_{1:t}^\top \mathbf{1}\mathbf{1}^\top X_{1:t} = (\mathbf{1}^\top X_{1:t})^2 = S_t^2$:
\[
X_{1:t}^\top\!\big(\Sigma_t^{(\sigma)}\big)^{\!-1}\! X_{1:t} \;=\; \frac{Q_t}{\sigma^2} - \frac{\tau_0^2 S_t^2/\sigma^4}{1 + t\tau_0^2/\sigma^2}.
\]
Assembling the three terms yields
\[
\log p(X_{1:t} \mid \sigma) \;=\; -\frac{t}{2}\log(2\pi) - \frac{t}{2}\log\sigma^2 - \frac{1}{2}\log\!\left(1 + \frac{t\tau_0^2}{\sigma^2}\right) - \frac{1}{2}\!\left[\frac{Q_t}{\sigma^2} - \frac{\tau_0^2 S_t^2/\sigma^4}{1 + t\tau_0^2/\sigma^2}\right],
\]
where the first term is absorbed into $\mathrm{const}(t)$ in equation~\eqref{eq:marginal-loglik}. The dependence on $X_{1:t}$ enters only through $(S_t, Q_t)$.

\bibliographystyle{plainnat}
\bibliography{refs}

\end{document}